\cotefile{157-1}
\scanpage{212}
\mspage{66}

Rappelons que les \deleted{flèches} \inserted{classes $\mathrm W_{\mathrm{asph}}$} \unsure{$W$-asphériques|homosphériques} et celles \inserted{$\mathrm{Liss}_{W}$} des flèches $W$-lisses, se déterminent mutuellement, ainsi

\[
\left\{\begin{aligned}
(f:X\to Y)\unsure{\in|\notin}\mathrm W_{\mathrm{asph}}&\Longleftrightarrow
  \forall\,(Y'\xrightarrow{u}Y)\in\mathrm{Liss}_{W},\\
  &\quad X'=X\times_Y Y'\longrightarrow X'\ \text{est}\ \in W\\
(\unsure{u|f}:Y'\to Y)\in\mathrm{Liss}_{W}&\Longleftrightarrow
  \text{le foncteur }X\mapsto X'=X\times_Y Y'\\
 &\quad\text{de }\mathrm{Cat}_{/Y}\text{ dans }\mathrm{Cat}_{/Y'}\ \text{transforme}\\
 &\quad\text{flèches dans }\mathrm W_{\mathrm{asph}}\text{ dans}\\
 &\quad\text{flèches dans }\mathrm W_{\mathrm{asph}}
\end{aligned}\right.
\tag*{(4.10)}
\]

Pour toute partie $A$ de $\mathrm{Fl}(\mathrm{Cat})$, \inserted{stable par isom.,} on peut associer une partie $\mathbb{L}_{\deleted{\unsure{W}}}(A)$ de $\mathrm{Fl}(\mathrm{Cat})$, par la relation

\[
\begin{aligned}
(u:Y'\to Y)\in\mathbb{L}_{\deleted{W}}(A)&\Longleftrightarrow u^*:\mathrm{Cat}_{/Y}\longrightarrow\mathrm{Cat}_{/Y'}\\
&\quad\text{transforme flèches}\\
&\quad\text{dans }A\text{ en flèches }\in A.
\end{aligned}
\tag*{(4.11.)}
\]

On voit que $\mathbb{L}_{\deleted{W}}(A)$ contient les iso, est stable par composition, \deleted{mais il n'est pas clair} \deleted{qu'il soit} \inserted{et est} stable aussi par chgt. de base (comme c'est le cas de $\mathrm{Liss}_{W}$). \deleted{\unsure{A d'|à d}} Inversement, à une partie $L$ de $\mathrm{Fl}(\mathrm{Cat})$, stable par isom., on associe une partie $\mathbb A_{\deleted{\unsure{W}}}(L)$ par

\[
\begin{aligned}
(f:X\to Y)\in\mathbb A_{\deleted{W}}(L)&\Longleftrightarrow\forall\,(Y'\to Y)\text{ dans }L,\\
&\quad X'\unsure{\Rightarrow|=}X\times_Y Y'\to Y'\\
&\quad\text{est dans }\deleted{\unsure{A|L}}W
\end{aligned}
\tag*{(4.12)}
\]

Si $L$ est stable par chgt. de base, alors $\mathbb A_{\deleted{W}}(L)$ est stable par composition. Il contient les iso \deleted{\unsure{si $L$ les}} (car $W$ les contient), \deleted{Mais il} \inserted{\unsure{est|et} contenu dans $W$,} \inserted{si $L$ contient les identités}. Mais
% archive sheet number: 212
\scanpage{213}
\mspage{67}

aucune raison qu'il soit stable par chgt de base quelconque (et en fait, $\deleted{\unsure{A}}\mathrm W_{\mathrm{asph}}$ ne l'est pas).

On a, pour une $A\subset W\subset\mathrm{Fl}(\mathrm{Cat})$ donnée \deleted{contenant} stable par iso,

\[
\begin{aligned}
 A\subset\mathbb A_W(\mathbb L(A))\overset{\mathrm{def}}{=}A^{\unsure{\sharp}}\quad
 \inserted{\subset W},\ &\text{partie stable}\\
 &\text{par composition,}\\
 &\text{contient iso}
\end{aligned}
\tag*{(4.13)}
\]

% slanted
\margin{\uline{\unsure{Canulé|Canular|annulé}}, \uline{cf. \S\unsure{5|6}}\par\rule{0.4pt}{11\baselineskip}}

et pour \unsure{tte|toute} partie $L$ de $\mathrm{Fl}(\mathrm{Cat})$ qui \deleted{\unsure{est}} stable par chgt de base, et par composition

\[
\begin{aligned}
 L\subset\mathbb L(\mathbb A_W(L))\overset{\mathrm{def}}{=}
 \unsure{L}^{\flat}\ &\text{partie de }\mathrm{Fl}\,\mathrm{Cat}\\
 &\text{stable par comp.}\\
 &\text{et par chgt. de base,}\\
 &\text{contient les iso}
\end{aligned}
\tag*{(4.14)}
\]

On a

\[
\left\{\begin{aligned}
 A^{\sharp\sharp}&=A^\sharp,\quad\text{car}\quad\mathbb A_W(L)^\sharp=\mathbb A_W(L)\\
 L^{\flat\flat}&=L^\flat,\quad\text{car}\quad\mathbb L(A)^\flat=\mathbb L(A)
\end{aligned}\right.
\]

Exemples :

\[
\begin{aligned}
1^\circ)\quad L&=\mathrm{Et},\ \mathbb A_W(\deleted{\unsure{E|Et}}L)=W_{\mathrm{asph}},\quad
  \mathbb L(\mathbb A_W(L))=L^\flat=\mathrm{Liss}_W\\
2^\circ)\quad L&=\unsure{\mathrm{Coet.}|\mathrm{Coét.}},\ \mathbb A_W(L)=W_{\unsure{\mathrm{coasp}|\mathrm{coasph}}},\quad
  \mathbb L(\mathbb A_W(L))=L^\flat=\mathrm{Prop}_W\\
3^\circ)\quad A&=W,\ \deleted{\mathbb A}\,\mathbb L(A)=\mathrm{Fib}_W,\quad A^\sharp=W\\
4^\circ)\quad L&=\mathrm{Fl}\,\mathrm{Cat},\ \mathbb A_W(L)=W_{\mathrm{univ}},\quad L^\flat=L.
\end{aligned}
\]
% archive sheet number: 213
\scanpage{214}
\mspage{69}

5. \uline{Dualité formelle}

Soient $P,Q$ deux ensembles ordonnés. \deleted{\unsure{Une|Nous}} J'\unsure{essaye|envisage|essaie} de dégager \deleted{la} \inserted{une} notion de ``dualité'' (formelle) \deleted{entre} \inserted{reliant} $P$ et $Q$. Ce sera une relation entre $P$ et $Q$

\[
R\subset P\times Q
\tag*{(5.1.)}
\]

ayant les propriétés suivantes

(D) Pour toute partie $A$ de $P$, l'ensemble

\[
\begin{aligned}
 A^{\unsure{\cup|\smile}}\overset{\mathrm{def}}{=}
 \big\{q\in Q\mid{}&\deleted{\unsure{R}}(p,q)\in R\ \deleted{\unsure{pour}\ \unsure{A}}\ \text{pour tout }p\in A\\
 &\text{i.e.}\quad A\times\{q\}\subset R\qquad\big\}
\end{aligned}
\]

est une partie ouverte de $Q$ ayant \inserted{un} plus grand élément $\beta(A)$\unsure{:|.} \inserted{$A^\cup=Q_{/\beta(A)}$.} Dualement, pour \deleted{toute} partie $B$ de $Q$, l'ensemble

\[
\begin{aligned}
 B^{\unsure{\cap|\frown}}\overset{\mathrm{def}}{=}
 \big\{p\in P\mid{}&(p,q)\in R\ \text{pour }\unsure{\text{tout}}\ q\in B,\\
 &\text{i.e.}\qquad \{p\}\times B\subset R\qquad\big\}
\end{aligned}
\]

est une partie ouverte de $P$ ayant un plus grand élément $\alpha(B)$, i.e. $B^\cap=P_{/\alpha(B)}$.

Le plus souvent, $P$ et $Q$ seront stables par sup quelconques (donc aussi par inf quelconques). Dans ce cas, (D) équivaut, du côté des $A^\cup$, (pour $A\in\unsure{\mathfrak P|\mathcal P}(P)$) à ceci

$(D_{0\unsure{a|\bullet|*}})$ Pour tout $p\in P$, $\deleted{\unsure{R(p)=\{q\in}}$ $R(p)=\{q\in Q\mid(p,q)\in R\}$ est une partie ouverte \uline{filtrante} de $\deleted{\unsure{P|R}}Q$, ou ce qui revient au $=$, une partie non vide ($\Longleftrightarrow$ contenant le plus petit élément $\varnothing_Q$ de $Q$) stable par sup de deux éléments. De plus $R(p)$ est stable par \deleted{\unsure{lim|les}} Sup filtrants, ou ce qui revient au $=$ (étant lui même filtrant), il contient un plus grand élément. Il revient
% archive sheet number: 214
\scanpage{215}
\mspage{70}

au même de dire que $R(p)$ est un \uline{ouvert} de $Q$ stable par Sup quelconques (ce qui implique qu'il contient $\varnothing_Q$). Dualement, pour tout $q\in Q$, $\,{}^{\unsure{s|t}}\!R(q)=\{p\in P\mid(p,q)\in R\}$ est une partie \uline{ouverte} de $P$, stable par Sup quelconques, i.e. une partie \unsure{ouverte|ouvert} \deleted{filtrante} ayant un plus grand élément.

En d'autres termes, quand on postule la condition (D) pour $A$ resp. $B$ réduit à un élément, alors \unsure{elle} reste vérifiée pour \unsure{tte|toute} partie. Cela provient des \deleted{\unsure{de|la}} formules

\[
\left\{\begin{aligned}
 \left(\bigcup_i A_i\right)^{\unsure{\cup|\smile}}&=\bigcap_i A_i^\cup\\
 \left(\bigcup_j B_j\right)^{\unsure{\cap|\frown}}&=\bigcap_j B_j^\cap,
\end{aligned}\right.
\tag*{(5.2.)}
\]

et du fait que l'intersection de parties ouvertes de la forme $P_{/\alpha_i}$ de $P$ est de la forme $P_{/\alpha}$, où $\alpha=\mathop{\mathrm{Inf}}_i\alpha_i$, et de $=$ dans $Q$. Ainsi on trouve

\[
\left\{\begin{aligned}
 \beta\left(\bigcup_i A_i\right)&=\mathop{\mathrm{Inf}}_i\beta(A_i)\\
 \alpha\left(\bigcup_j B_j\right)&=\mathop{\mathrm{Inf}}_j\alpha(B_j).
\end{aligned}\right.
\tag*{(5.3.)}
\]

On posera, pour une partie $A$ de $P$ resp. $B$ de $Q$

\[
 A^{\unsure{\sharp}}=(A^\cup)^\cap,\qquad B^\flat=(B^\cap)^\cup
\tag*{(5.4).}
\]

\deleted{\unsure{(5.5)}} de sorte qu'on aura

\[
 A\subset A^\sharp,\qquad B\subset B^\flat
\tag*{(5.5.)}
\]

\deleted{On aura, pour tt $A$ resp.}

\deleted{\[
 (A^\cup)
 \tag*{(5.6)}
\]}

On dira que $A$ est \uline{$R$-clos} si $A=A^\sharp$ \inserted{(ou encore si $A$ est de la forme $B'^\cap$ pour $B'\subset Q$, cf 5.6)}, que
% archive sheet number: 215
\scanpage{216}
\mspage{\deleted{\unsure{7|8}}71\quad A}

$B$ est \uline{$R$-clos si} $B=B^\flat$. On aura pour tout \deleted{$A$} $A\subset P$, $B\subset Q$

\[
\left\{\begin{aligned}
 (A^{\unsure{\cup|\smile}})^\flat&=A^\cup\quad\text{i.e. }A^\cup\text{ est }R\text{-clos}\\
 (B^{\unsure{\cap|\frown}})^{\unsure{\sharp}}&=B^\cap\quad\text{i.e. }B^\cap\text{ est }R\text{-clos.}
\end{aligned}\right.
\tag*{(5.6)}
\]

En particulier

\[
 (A^\sharp)^\sharp=A^\sharp,\qquad (B^\flat)^\flat=B^\flat
\tag*{(5.7.)}
\]

i.e. la $R$-\uline{clôture} $A^\sharp$ de $A$ \inserted{$\subset P$,} est $R$-close, et de $=$ la $R$-clôture $B^\flat$ de $B\subset Q$.

En fait, je m'aperçois que je suis en train de mélanger deux espèces de faits différents. \deleted{\unsure{D'}} Les relations (5.2.), et (\unsure{5.4.|5.3.}) à (5.7), ne dépendent pas de la donnée des relations d'ordre sur $P,Q$, ni d'axiomes reliant celles-ci à la relation donnée $R$, mais seulement de cette relation, laquelle n'a pas à être soumise à aucune condition. Notons que

\[
 \varnothing^\cup=\deleted{\unsure{B}}Q,\qquad \varnothing^\cap=\deleted{\unsure{A}}P
 \tag*{(5.8.)}
\]

tandis que $\deleted{\unsure{\varnothing^\sharp=}}$

\[
 \varnothing^\sharp=\unsure{Q}^\cap,\qquad \varnothing^\flat=P^\cup
 \tag*{(5.9)}
\]

sont respectivement la plus petite partie $\deleted{\unsure{cl}}R$-close de $P$ et de $Q$.

L'application

\[
 A\unsure{\vdash|\mapsto} A^\sharp
 \tag*{(5.10\unsure{.|:})}
\]

est une application \uline{croissante} \inserted{(et \uline{idempotente})} de $\deleted{\unsure{\mathfrak P A}}$ \deleted{dans} $\unsure{\mathfrak P|\mathcal P}(P)$ dans lui-même, ayant \inserted{\unsure{donc|aussi}} la propriété

\deleted{\[
 \unsure{(A\subset B)^\sharp=A}
\]}
% archive sheet number: 216
\scanpage{217}
\mspage{72}

\[
 \left(\bigcup A_i\right)^{\unsure{\sharp}}=\left(\bigcup_i A_i^\sharp\right)^\sharp\ ;
\]

Pour tout $A\inserted{\subset P}$, $A^\sharp$ est la plus petite partie $R$-close de $P$ contenant $A$, et de $=$ pour tt $B\subset Q$, $B^\flat$ est la plus petite partie $R$-close de $Q$ contenant $B$. On a de plus (c'était $=$ à dire avant de dire que $A\mapsto A^\sharp$ est croissante)

\[
\left\{\begin{aligned}
 A\subset A'&\Longrightarrow A'^{\unsure{\cup|\smile}}\subset A^\cup\\
 B\subset B'&\Longrightarrow B'^{\unsure{\cap|\frown}}\subset B^\cap,
\end{aligned}\right.
\tag*{(5.11)}
\]

i.e. les applications

\diagram{\begin{tikzcd}[ampersand replacement=\&,column sep=large]
\text{(5.12)} \& \unsure{\mathfrak P|\mathcal P}(\unsure{p|P})\arrow[r,shift left,"\unsure{B|\beta}"]
 \& \mathfrak P(Q)\arrow[l,shift left,"\alpha"]
 \& \begin{aligned}A&\unsure{\vdash|\mapsto} A^\cup\\ B&\vdash B^\cap\end{aligned}
\end{tikzcd}}

sont décroissantes \unsure{(ce|Ce} qui implique que les deux composés $\alpha\beta=(A\mapsto A^\sharp)$ et $\beta\alpha=(B\mapsto B^\flat)$ sont croissantes). Leurs restrictions aux sous-ensembles $\mathfrak P_R(P)$ et $\mathfrak P_{\unsure{{}^{s}R|{}^{t}R}}(Q)$ des parties $R$-closes de $P$ resp. de $Q$, sont \deleted{\unsure{donc} telles}, par définition de celles-ci comme

\[
\begin{aligned}
 \mathfrak P_R(P)&=\{A\subset P\mid A^\sharp(\overset{\mathrm{def}}{=}\alpha\beta(A))=A\}\\
 \mathfrak P_R(Q)&=\{B\subset Q\mid B^\flat(\overset{\mathrm{def}}{=}\beta\alpha(\unsure{B,|B}))=B\},
\end{aligned}
\]

inverses l'une de l'autre. On trouve ainsi

\uline{Proposition} 5.1\quad Soient $P,Q$ deux ensembles, $R\inserted{\subset P\times Q}$ une relation. Alors les applications $\alpha,\beta$ de (5.12.) \unsure{induisent} des bijections, inverses l' 
% archive sheet number: 217
\scanpage{218}
\mspage{73}

une de l'autre, entre l'ensemble $\mathrm{Cl}_R(P)$ des parties $R$-closes de $P$, et l'ensemble $\mathrm{Cl}_{\unsure{{}^{s}R|{}^{t}R}}(Q)$ des parties $\,{}^{s}R$-closes (nous dirons $R$-closes par abus de langage) de $Q$. Ces bijections sont \inserted{décroissantes, donc} des anti-\deleted{\unsure{in|t}}isomorphismes pour les relations d'ordre naturelles (d'inclusion).

\margin{\rule{0.4pt}{4\baselineskip}}%
\uline{Corollaire} 5.2. Soit $(A_i)_i$ une famille de parties \uline{$R$-closes} de $P$. Alors $\bigcap A_i$ est une partie $R$-close de $P$, et

\[
 \left(\bigcap A_i\right)^{\unsure{\cup|\smile}}=
 \underbrace{\left(\bigcup_i A_i^\cup\right)^\flat}_{
 \substack{\text{\deleted{plus petite} Sup des}\\\text{parties }R\text{-closes}\\A_i^\cup\text{ de }Q}}.
 \tag*{(5.2.1)}
\]

Mais la relation $R$ \unsure{se} ``prolonge'' en une relation $\overline R$ entre

\[
 \overline P=\unsure{\mathfrak P|\mathcal P}(P),\qquad\overline Q=\mathfrak P(Q)
 \tag*{(5.13)}
\]

par la prescription

\[
\begin{aligned}
 \deleted{(A,B)\in}\underset{\substack{\cap\\\overline P\times\overline Q}}{\overline R}
  =\big\{(A,B)\in\overline P\times\overline Q\mid{}&
  B\subset A^\cup\quad\text{i.e. }A\subset B^{\unsure{\cap|\frown}}\\
 &\text{i.e. }A\times B\subset R\big\}
\end{aligned}
\tag*{(5.14)}
\]

\deleted{\unsure{Cette}} Les ensembles $\overline P,\overline Q$ sont stables par Sup quelconques (\unsure{réunions|réunion} de parties de $P$, resp $Q$), et la relation $\overline R$ sur ces ensembles ordonnés satisfait à la condition $(D_0)$
% archive sheet number: 218
\scanpage{219}
\mspage{74}

ci-dessus (p. 69), de façon évidente : pour $A$ fixé, les $B$ tels que $(A,B)\in\overline R$ sont ceux qui sont contenus dans $A^{\unsure{\cup|\smile}}$, donc l'ensemble $\overline{\unsure{R|Q}}(A)$ n'est autre que $\overline Q_{/A^\cup}$, et dualement $\,{}^{\unsure{s|t}}\overline R(B)$ n'est autre que $\overline P_{/B^{\unsure{\cap|\frown}}}$. Ainsi

\[
 \beta(A)=A^\cup,\qquad\alpha(B)=B^\cap.
 \tag*{(5.14)}
\]

Si j'ai songé à introduire \deleted{\unsure{dans}} l'axiome $D_0$, pour \deleted{des} \inserted{un couple} d'ensembles ordonnés quelconques, c'est à cause des cas où ceux-ci ne sont \uline{pas} de la forme $\unsure{\mathfrak P|\mathcal P}(P_{\unsure{0|\sharp|,}}),\mathfrak P(Q_{\unsure{,|0}})$, et (alors \unsure{même|mêmes} qu'ils le seraient) où la relation $\unsure{R}$ \inserted{de base} n'est pas déduite \inserted{comme ci-dessus} d'une relation quelconque entre ensembles $P,Q$.

Je reviens \unsure{donc|sur} à l'axiomatique du début, avec $P,Q$ ordonnés, et $R$ satisfaisant (D). Dans tous les cas auxquels je peux songer, $P$ et $Q$ sont stables par Sup quelconques, donc (D) équivaut à $(D_0)$.

Notons que l'application

\[
 P\longrightarrow\mathfrak P(P),\qquad p\unsure{\vdash|\mapsto} P_{/p}
\]

est injective, et de $=$ pour $Q\to\mathfrak P(Q)$. Ainsi, la connaissance de $A\mapsto A^\cup$ équivaut à celle de $A\longmapsto\beta(A)$

\[
 \beta:\mathfrak P(P)\longrightarrow Q,
\]

et de $=$ celle de $B\vdash B^\cap$ équivaut à
% archive sheet number: 219
\scanpage{220}
\mspage{75}

celle de

\[
 \alpha:\unsure{\mathfrak P|\mathcal P}(Q)\longrightarrow P.
\]

\deleted{La relation $R$ se retrouve à partir \unsure{de|des}} En \unsure{fait|fais}, il suffit de se donner

\[
 \left\{\begin{aligned}
  \beta_0 &:P\longrightarrow Q\\
  \alpha_0 &:Q\longrightarrow P
 \end{aligned}\right.\qquad
 \begin{aligned}&\text{applications}\\&\text{\uline{décroissantes}}\end{aligned}
 \tag*{(5.15)}
\]

\deleted{\unsure{q}} définies par

\[
 \left\{\begin{aligned}
 \beta_0(p)&=\beta(\{p\})\\
 \alpha_0(q)&=\alpha(\{q\}),
 \end{aligned}\right.
\]

puisqu'on en déduit $\beta,\alpha$ par

\[
 \left\{\begin{aligned}
 \beta(A)&=\mathop{\mathrm{Inf}}_{p\in A}\beta(p)\\
 \alpha(B)&=\mathop{\mathrm{Inf}}_{q\in B}\alpha(q).
 \end{aligned}\right.
 \tag*{(5.16)}
\]

La relation $R$ se déduit de la donnée de $\beta_0$, ou du choix de $\alpha_0$, par

\[
 (p,q)\in R\Longleftrightarrow q\leq\beta_0(p)\Longleftrightarrow p\leq\alpha_0(q).
 \tag*{(5.17)\unsure{\textsubscript{e}|\textsubscript{a}|.}}
\]

Cela implique donc que $\alpha_0,\beta_0$ se déterminent mutuellement, par les formules

\[
 \left\{\begin{aligned}
 \alpha_0(q)&=\deleted{\unsure{\mathrm{Inf}}}\mathrm{Sup}\,
 \underbrace{\{p\in P\mid q\leq\beta_0(p)\}}_{\{q\}^{\unsure{\cap|\frown}}}\\
 \beta_0(p)&=\mathrm{Sup}\,
 \underbrace{\{q\in Q\mid p\leq\alpha_0(q)\}}_{\{p\}^{\unsure{\cup|\smile}}}.
 \end{aligned}\right.
 \tag*{(5.18)}
\]

Reste à se demander quelles sont les applications $\alpha_0:P\unsure{-|\to}Q$, ou $\beta_0:Q\unsure{\to|-}P$, qui correspondent à une \uline{``relation dualisante''}
% archive sheet number: 220
\scanpage{221}
\mspage{76}

$R$ sur $P\times Q$. Voici les conditions sur $\beta_0:P\to Q$ :

\textcircled{a}\quad $\beta_0$ est \uline{\unsure{décroissante|décroissant}},

\deleted{b) \unsure{pour} $\forall$ partie} \unsure{ce} qui implique que la relation

\[
 \underset{\substack{\cap\\P\times Q}}{R_{\beta_0}}\deleted{\subset}
 =\{(p,q)\in P\times Q\mid q\leq\beta_0(p)\}
\]

satisfait à la condition $(D_0)$ du côté gauche, et il faut encore la vérifier du côté droit, i.e regarder, pour $q$ fixé, l'ensemble

\[
 \{q\}^{\unsure{\cap|\frown}}=\{p\in P\mid q\leq\beta_0(p)\}.
\]

Comme $\beta_0$ est \unsure{décroissante|décroissant}, on voit que \deleted{\unsure{est}} $p\in\{q\}^\cap$, et $p'\leq p\Longrightarrow p'\in\{q\}^\cap$, i.e. $\{q\}^\cap$ est \uline{ouvert}. \deleted{Pour} On doit vérifier encore (supposant l'existence des Sup. quelconques dans $P$) qu'il \deleted{est} cet ensemble est stable par \deleted{\unsure{lim}} Sup. quelconques - il en résultera qu'il est de la forme $P_{/\alpha(q)}$. Cela équivaut à la condition

\textcircled{b}\quad $\beta_0$ transforme Sup quelconques \inserted{dans $P$} en Inf \inserted{dans $Q$}.

\deleted{\unsure{(}} Si on regarde $\beta_0$ comme un \inserted{\unsure{contr.|covar.}} foncteur $P^{\unsure{\circ|o}}\longrightarrow Q$, cela signifie qu'il commute aux limites projectives, ou encore qu'il a un adjoint à gauche $Q\longrightarrow P^\circ$, lequel commute aux limites inductives, et correspond
% archive sheet number: 221
\scanpage{222}
\mspage{77}

à l'application décroissante $\alpha_0:Q\unsure{-|\to}P$, transformant Sup en Inf. \deleted{L'application}

Ici il \deleted{peut} est commode de considérer $\unsure{\sharp}$ et $\flat$ comme des opérations sur $P$ resp $Q$, par

\[
 p^\sharp=\alpha_0\beta_0(p),\qquad q^\flat=\beta_0\alpha_0(q).
 \tag*{(5.19)}
\]

La relation entre ces endomorphismes, et l'opération de $=$ nom sur $\unsure{\mathfrak P|\mathcal P}(\deleted{\unsure{A}}P),\mathfrak P(Q)$, est donnée par

\[
 \left\{\begin{aligned}
 p^\sharp&=\text{\unsure{Plus|plus} grand élément de }\{p\}^\sharp\\
 A^\sharp&=\deleted{\unsure{p}}\mathop{\mathrm{Sup}}_{p\in A}p^\sharp,
 \end{aligned}\right.
 \tag*{(5.20)}
\]

et symétriquement du côté de $Q$.

\deleted{Les applications} On dit que $p$ est \uline{$R$-clos}, ou $\beta_0$-clos, ou $\alpha_0$-clos, si $p=p^\sharp$, et de $=$ $q\in Q$ est dit $R$-\uline{clos} etc si $q=q^\flat$. Les \deleted{appl} applications

\[
 \left\{\begin{aligned}
 p\unsure{\vdash|\mapsto} p^\sharp &: P\longrightarrow P\\
 q\vdash q^\flat &: Q\unsure{-|\to}Q
 \end{aligned}\right.
 \tag*{(5.21)}
\]

sont \uline{croissantes}, \inserted{et \uline{idempotentes},} \deleted{et commutent aux Sup \uline{quelconques}. Ce dernier fait résulte de} \inserted{on peut caractériser} $p^\sharp$ comme le plus petit élément clos de $P$ majorant $p$, et de $=$ pour $q^\flat$. On trouve \unsure{alors}, par (5.19), que $\alpha_0$ et $\beta_0$ induisent des bijections inverses l'une de l'autre,
% archive sheet number: 222
\scanpage{223}
\mspage{78}

entre l'\unsure{ens.} des \unsure{élément|éléments} clos de $P$, et l'ens des \unsure{éléments} clos de $Q$

\diagram{\begin{tikzcd}[ampersand replacement=\&,column sep=large]
\text{(5.22.)} \& \unsure{p|P}^{\unsure{\sharp}}
 \arrow[r,shift left,"{\substack{\overline{\beta}_0\\\sim}}"] \&
 Q^\flat\arrow[l,shift left,"\overline{\alpha}_0","\sim"'] .
\end{tikzcd}}

Mais en pratique, les applications $\alpha_0,\beta_0$ ne sont pas obtenues directement, mais on dispose de la relation $R$, satisfaisant la condition $(D_0)$, dont on déduit $\alpha_0,\beta_0$. Dans tous les cas auxquels je \deleted{\unsure{connais}} \inserted{pense} ici, ou bien $P,Q$ \unsure{sont} de la forme $\unsure{\mathfrak P|\mathcal P}(P_0),\mathfrak P(Q_0)$, et $R$ se déduit d'une relation $R_0\subset P_0\times Q_0$ comme ci-dessus (cas de loin le plus fréquent, il me semble), ou bien $P$ et $Q$ apparaissent comme sous-ensembles \inserted{convenables} d'ensembles $\mathfrak P(P_0),\mathfrak P(Q_0)$, munis de l'ordre induit, et la relation $R$ sur $P\times Q$ s'obtient d'une façon plus complexe.

\uline{Exemples}

1) Relation d'\uline{orthogonalité} entre \deleted{parties} de deux modules sur un anneau, \deleted{\unsure{munis d'un}} d'un accouplement de dualité. \inserted{Dans le cas de \unsure{vectoriels|modules} de dimension finie et d'une vraie dualité, les parties closes sont les sous-esp. vectoriels.}

% margin note written sideways
\margin{\deleted{\unsure{dans, ce sont}} dans \unsure{le|ce} cas \unsure{infini|en fait}, \unsure{se|ce} sont les sous-espaces vectoriels fermés pour la ``\unsure{topologie|topol.} \unsure{faible}'' définie par l'accouplement.}

2) Relations d'\uline{orthogonalité} entre parties de deux groupes topologiques \unsure{comm.} loc. \unsure{cpcts|cpacts} duaux l'un de l'autre, au sens de \unsure{Pontragin|Pontriagin|Pontrjagin|Pontraguin}.
% archive sheet number: 223
\scanpage{224}
\mspage{79}

Les parties closes sont les sous-groupes fermés.

3) Pour deux espaces vectoriels ``\unsure{en} dualité'' sur un corps valué complet, tel $\mathbb R$ ou $\mathbb C$, la relation

\[
 |\langle x,x'\rangle|\leq 1.
\]

Dans le cas réel ou complexe, les parties closes sont les parties convexes \unsure{cerclées} faiblement fermées

4) Pour les espaces vectoriels réels, la relation de semi-polarité

\[
 \langle x,x'\rangle\deleted{\unsure{>|\geq}}\unsure{\geq|\geqq}-1.
\]

Les parties closes sont les parties convexes \unsure{fblt.} fermées contenant l'origine.

5) Soit $K$ un corps, \inserted{(pour simplifier)} $L$ une extension galoisienne de groupe $G$, Relation entre $L$ et $G$

\[
 \sigma.x=x\qquad\text{i.e. }x\in L^\sigma.
\]

Les parties closes de $L$ sont les sous-extensions, les parties closes de $G$ les sous-groupes de $G$

Il y a des \deleted{\unsure{des|des deg.}} analogues divers \unsure{dans des|dans les|de ces} cas \unsure{inséparables} (à \unsure{coup\deleted{t}s|coups|compte} de $p$-algèbres de Lie par exemple), et des énoncés de descente en tous genres.
% archive sheet number: 224
\scanpage{225}
\mspage{80}

\deleted{\unsure{6)}} Tous ces exemples sont de $1^{\mathrm{ere}}$ espèce, les ensembles ordonnés qui interviennent sont de la forme $\unsure{\mathfrak P|\mathcal P}(P_0),\mathfrak P(Q_0)$, et la relation entre eux est déduite d'une relation $R_0$ entre $P_0$ et $Q_0$. \unsure{voici|Voici} maintenant des exemples plus complexes.

6) Soit $M$ une catégorie (``grosse'', le plus souvent \ldots), On prend $\unsure{\mathfrak P|P}$

\[
 \unsure{\mathfrak P|P}=Q=\mathfrak P(\mathrm{Fl}(M)).
\]

La relation $R\subset P\times Q$ est la relation de Quillen, que j'écris, pour $\Phi,\Psi\in\mathfrak P$

\[
 \begin{aligned}
 &\deleted{\unsure{B\subset A,}}\\
 &\deleted{\unsure{\Phi}}\Psi\subset\Phi_*\,\unsure{,|?}\quad\deleted{\unsure{Q|R}}
 \end{aligned}
\]

ou encore de façon équivalente

\[
 \Phi\subset\Psi^*
\]

($\unsure{\Psi|\Phi|\Psi^{+}}$ a la LLP pour $\Psi$, $\Phi$ la RLP pour $\Psi$). \deleted{\unsure{Les}} \deleted{Donc} les applications $\alpha_0,\beta_0$ sont données par

\[
 \beta_0(\Phi)=\Phi_*,\qquad \alpha_0(\Psi)=\Psi^*.
\]

L'idée ici ne viendrait à personne d'introduire $\mathfrak P(P),\mathfrak P(Q)$ \deleted{\unsure{et}} c'est déjà assez gros comme ça ! La notation pour $\Phi^{\unsure{\sharp}}$ est ici $\unsure{\widetilde\Phi|\overline\Phi}$, pour $\unsure{\widehat{\Psi}|\widetilde{\Psi}|\Psi}^\flat$ c'est $\overline\Psi$. Les parties closes sont celles que j'appelle $Q$-closes à g. et à droite.

La relation $R$, que j'écris

\[
 \Phi\longleftrightarrow\Psi,
\]
% archive sheet number: 225
\scanpage{226}
\mspage{81}

n'est pas déduite d'une relation entre flèches de $M$ i.e d'une partie de $\mathrm{Fl}(M)\times\mathrm{Fl}(M)$, mais est de nature plus complexe. Pour la décrire, il faut introduire l'ensemble $D$ des diagrammes \inserted{carrés comm. (traits pleins)} de la forme

\diagram{\begin{tikzcd}[ampersand replacement=\&]
{}\arrow[d,phantom,"{(\unsure{\alpha})}" description] \& A\arrow[r]\arrow[d,"i"] \& X\arrow[d,"p"]\\
 {} \& B\arrow[r]\arrow[ur,dashed,"{\unsure{f|g}}"] \& Y
\end{tikzcd}}

(qu'on peut interpréter comme $\mathrm{Fl}(\underline{\mathrm{Fl}}(M))$, ensemble des flèches dans $\underline{\mathrm{Fl}}(M)$)\,;

On a une partie $D_0$ de $D$, \unsure{formée} des diagrammes pour lesquels existe $f$ avec deux triangles commutatifs (c'est donc l'image de $\underset{\mathrm{Hom}(\unsure{\mathbb{\Delta}|\Delta}^{3},M)}{\underline{\mathrm{Fl}_3(M)}}$ dans $\mathrm{Hom}(\mathbb{\Delta}^1\times\mathbb{\Delta}^1,M)$, transposée de $\mathbb{\Delta}^1\times\mathbb{\Delta}^1\longrightarrow\mathbb{\Delta}^3$ donné par $(0,0)\mapsto0$, $(0,1)\mapsto1$, $(1,0)\unsure{\vdash|\mapsto}2$, $(1,1)\vdash3$), et une application

\diagram{\begin{tikzcd}[ampersand replacement=\&,row sep=1.1em]
 D\arrow[r,"{\unsure{\varphi|\phi}}"] \& \mathrm{Fl}(M)\times\mathrm{Fl}\unsure{(}M),\\
 D_0\arrow[u,phantom,"\cup" description] \& {}
\end{tikzcd}}

associant à $\alpha$ le couple $(i,\unsure{j|p})$ (source et but de $\alpha$, considéré comme \unsure{flèche|flèches} dans $\underline{\mathrm{Fl}}(M)$ de $i$ dans $p$). La relation $R$ est alors déduite de $\varphi,D_0$ par la condition
% archive sheet number: 226
